\documentclass[10pt]{article}
\usepackage[utf8]{inputenc}
\usepackage[T1]{fontenc}
\usepackage{amsmath}
\usepackage{amsfonts}
\usepackage{amssymb}
\usepackage[version=4]{mhchem}
\usepackage{stmaryrd}
\usepackage{graphicx}
\usepackage[export]{adjustbox}
\graphicspath{ {./images/} }
\usepackage{bbold}

\begin{document}
\section*{Bilateral trade}
\section*{Model: Symmetric IPVM plus regularity (i.e, increasing $v_{i}(\cdot)$ )}
\begin{itemize}
  \item Classic auctions with appropriate reserve price maximize seller's expected revenue.
  \item When maximizing expected seller's revenue, it is fitting to consider the seller as the mechanism designer and to assume that the seller knows her valuation.
  \item Classic autions without reserve price maximize social surplus provided the seller's valuation is known.
\end{itemize}

\section*{Model: IPVM without symmetry or regularity}
\begin{itemize}
  \item Without symmetry or regularity results are "fewer."
  \item The revenue maximizing DRM is known (Myerson (1981)). It is not obvious how to use it. The economic implications are not as developed.
\end{itemize}

\section*{Market design, efficiency of bilateral trade}
\begin{itemize}
  \item Are trading mechanisms "efficient" in the sense that trade takes place always if the buyer's valuation exceeds the seller's valuation and never otherwise?
  \item Appropriate model:
  \item Seller and buyer have private information (types);
  \item the mechanism designer does not observe agents' types.
  \item IPVM variation: a buyer and a seller with private information, $x_{b}, x_{s}$ respectively, r.v. in $[c, d]$.
  \item Question. Is there an ex post efficient mechanism? That is, is there a mechanism that ensures that trade takes place if and only if $x_{b}>x_{s}$ ?
\end{itemize}

\section*{Fixed uninformed agent's valuation}
\includegraphics[max width=\textwidth, alt={}, center]{39af411a-b5cd-4820-bf76-1645c3db63ba-05_1048_820_348_154}\\
\includegraphics[max width=\textwidth, alt={}, center]{39af411a-b5cd-4820-bf76-1645c3db63ba-05_1050_815_348_1034}\\
\includegraphics[max width=\textwidth, alt={}, center]{39af411a-b5cd-4820-bf76-1645c3db63ba-05_1045_815_351_1915}

\section*{Bilateral private information}
\begin{itemize}
  \item Both buyer and seller observe their valuations $x_{b}$ and $x_{s}$ distributed according to pdfs following $f_{b}\left(x_{b}\right)$ and $f_{s}\left(x_{s}\right)$ respectively.
  \item See Vickrey (1961). There is no mechanism with the following properties:\\
i. truthtelling is a DS equilibrium;\\
ii. no outside subsidies;\\
iii. ex post Pareto efficiency.
  \item Both d'Aspremont and Gérard-Varet (1979) and Arrow (1979) find mechanism with\\
i. Bayesian IC;\\
ii. no outside subsidies but no IR constraints,\\
iii. ex post efficiency.
  \item Myerson and Satterthwaite (1983) prove that no mechanism satisfies (i) Bayesian IC; (ii) IR and no outside subsidies; (iii) ex post efficiency.
\end{itemize}

\section*{Bilateral trade-Theorem}
Theorem 1 (Myerson and Satterthwaite) IPVM with 2 agents, a buyer and a seller, with valuations $x_{b} \in[c, d]$ and $x_{s} \in[c, d]$ respectively. (Valuations are independently distributed, not necessarily identically.)

There is no mechanism that is (i) Bayesian IC; (ii) satisfies IR (without outside resources); and (iii) is efficient.

\begin{itemize}
  \item The supports need to overlap but do not need to coincide.
\end{itemize}

\section*{Proof.}
\begin{itemize}
  \item Let $(\tilde{p}, \tilde{t})$ be a BIC mechanism where
\end{itemize}

$$
(\tilde{p}, \tilde{t})=\left(p_{b}\left(x_{b}, x_{s}\right), t_{b}\left(x_{b}, x_{s}\right), p_{s}\left(x_{b}, x_{s}\right), t_{s}\left(x_{b}, x_{s}\right)\right)
$$

with $p_{b}\left(x_{b}, x_{s}\right)+p_{s}\left(x_{b}, x_{s}\right)=1$ and $t_{b}\left(x_{b}, x_{s}\right)+t_{s}\left(x_{b}, x_{s}\right)=0$.

\begin{itemize}
  \item ( $\tilde{p}, \tilde{t}$ ) may be described by two functions instead of four,
  \item $p:[c, d]^{2} \rightarrow[0,1]$, and $t:[c, d]^{2} \rightarrow \mathbb{R}$ where
  \item $p_{b}=p$ (probability that buyer gets the object),
  \item $p_{s}=1-p_{b}$ (probability seller keeps object),
  \item $t_{b}=t$ (payment from buyer to seller),\\
$t_{s}=-t_{b}$ (amount received by seller).
\end{itemize}

Then,

$$
\begin{gathered}
u_{b}\left(x_{b}, x_{s}\right)=p\left(x_{b}, x_{s}\right) x_{b}-t\left(x_{b}, x_{s}\right) \\
u_{s}\left(x_{b}, x_{s}\right)=t\left(x_{b}, x_{s}\right)-p\left(x_{b}, x_{s}\right) x_{s} \\
U_{b}\left(x_{b}\right)=E_{x_{s}}\left[u_{b}\left(x_{b}, x_{s}\right)\right]=Q_{b}\left(x_{b}\right) x_{b}-T_{b}\left(x_{b}\right) \\
U_{s}\left(x_{s}\right)=E_{x_{b}}\left[u_{s}\left(x_{b}, x_{s}\right)\right]=T_{s}\left(x_{s}\right)-Q_{s}\left(x_{s}\right) x_{s}
\end{gathered}
$$

where

$$
\begin{aligned}
T_{s}\left(x_{s}\right) & =E_{x_{b}} t\left(x_{b}, x_{s}\right) \\
Q_{s}\left(x_{s}\right) & =E_{x_{b}} p\left(x_{b}, x_{s}\right)
\end{aligned}
$$

The mechanism must satisfy two constraints

\begin{itemize}
  \item Interim individual rationality\\
$\left(\forall x_{b}\right) U_{b}\left(x_{b}\right) \geq 0$ and $\left(\forall x_{s}\right) U_{s}\left(x_{s}\right) \geq 0$.
  \item Bayesian incentive compatibility (BIC),
\end{itemize}


\begin{align*}
& U_{s}\left(x_{s}\right)=U_{s}(d)+\int_{x_{s}}^{d} Q_{s}(y) \mathrm{d} y  \tag{1}\\
& U_{b}\left(x_{b}\right)=U_{b}(c)+\int_{c}^{x_{b}} Q_{b}(y) \mathrm{d} y
\end{align*}


When players observed $\left(x_{b}, x_{s}\right)$, total interim surplus

$$
U_{s}\left(x_{s}\right)+U_{b}\left(x_{b}\right)=Q_{b}\left(x_{b}\right) x_{b}-T_{b}\left(x_{b}\right)+T_{s}\left(x_{s}\right)-Q_{s}\left(x_{s}\right) x_{s} .
$$

Since $-E_{x_{b}} T_{b}\left(x_{b}\right)+E_{x_{s}} T_{s}\left(x_{s}\right)=0$, expected total surplus is


\begin{align*}
& E\left[U_{s}\left(x_{s}\right)+U_{b}\left(x_{b}\right)\right]=E\left[Q_{b}\left(x_{b}\right) x_{b}-Q_{s}\left(x_{s}\right) x_{s}\right] \\
& =\int_{c}^{d} \int_{c}^{d}\left[Q_{b}\left(x_{b}\right) x_{b}-Q_{s}\left(x_{s}\right) x_{s}\right] f_{b}\left(x_{b}\right) f_{s}\left(x_{s}\right) \mathrm{d} x_{b} \mathrm{~d} x_{s} \\
& E\left[U_{s}\left(x_{s}\right)+U_{b}\left(x_{b}\right)\right] \\
& \quad=\int_{c}^{d} \int_{c}^{d} p\left(x_{b}, x_{s}\right)\left[x_{b}-x_{s}\right] f_{b}\left(x_{b}\right) f_{s}\left(x_{s}\right) \mathrm{d} x_{b} \mathrm{~d} x_{s} \tag{2}
\end{align*}


From Equation 1 (IC, integrability condition)

$$
\begin{aligned}
E\left[U_{s}\left(x_{s}\right)\right. & \left.+U_{b}\left(x_{b}\right)\right]=U_{s}(d) \\
& +\int_{c}^{d}\left[\int_{x_{s}}^{d} Q_{s}(y) \mathrm{d} y\right] f_{s}\left(x_{s}\right) \mathrm{d} x_{s} \\
& +U_{b}(c)+\int_{c}^{d}\left[\int_{c}^{x_{b}} Q_{b}(z) \mathrm{d} z\right] f_{b}\left(x_{b}\right) \mathrm{d} x_{b}
\end{aligned}
$$

Integrating by parts both terms,

$$
\begin{aligned}
\int_{c}^{d} & {\left[\int_{c}^{x_{b}} Q_{b}(z) \mathrm{d} z\right] f_{b}\left(x_{b}\right) \mathrm{d} x_{b} } \\
& =\int_{c}^{d} Q_{b}\left(x_{b}\right) \mathrm{d} x_{b}-\int_{c}^{d} F_{b}\left(x_{b}\right) Q_{b}\left(x_{b}\right) \mathrm{d} x_{b} \\
& =\int_{c}^{d}\left[1-F_{b}\left(x_{b}\right)\right] Q_{b}\left(x_{b}\right) \mathrm{d} x_{b}
\end{aligned}
$$

and

$$
\int_{c}^{d}\left[\int_{x_{s}}^{d} Q_{s}(y) \mathrm{d} y\right] f_{s}\left(x_{s}\right) \mathrm{d} x_{s}=\int_{c}^{d} F_{s}\left(x_{s}\right) Q_{s}\left(x_{s}\right) \mathrm{d} x_{s}
$$

$$
\begin{aligned}
& E\left[U_{s}\left(x_{s}\right)+U_{b}\left(x_{b}\right)\right]=U_{s}(d)+U_{b}(c)+ \\
& \int_{c}^{d} \int_{c}^{d}\left\{F_{s}\left(x_{s}\right) f_{b}\left(x_{b}\right)+\left[1-F_{b}\left(x_{b}\right)\right] f\left(x_{s}\right)\right\} p\left(x_{b}, x_{s}\right) \mathrm{d} x_{b} \mathrm{~d} x_{s}
\end{aligned}
$$

Equating the two expressions of total surplus, Equation 3 and Equation 2, and solving, we obtain

$$
\begin{aligned}
& U_{s}(d)+U_{b}(c)= \\
& =\int_{c}^{d} \int_{c}^{d}\left\{\left[x_{b}-\frac{1-F_{b}\left(x_{b}\right)}{f_{b}\left(x_{b}\right)}\right]-\left[x_{s}+\frac{F_{s}\left(x_{s}\right)}{f_{s}\left(x_{s}\right)}\right]\right\} \\
& \quad p\left(x_{b}, x_{s}\right) f_{s} f_{b} \mathrm{~d} x_{b} \mathrm{~d} x_{s} \\
& \geq 0
\end{aligned}
$$

Where the inequality follows by IR. Suppose

$$
p\left(x_{b}, x_{s}\right)= \begin{cases}1 & \text { if } x_{b} \geq x_{s} \\ 0 & \text { otherwise }\end{cases}
$$

$$
\begin{aligned}
& \int_{c}^{d} \int_{c}^{d}\left\{\left[x_{b}-\frac{1-F_{b}\left(x_{b}\right)}{f_{b}\left(x_{b}\right)}\right]-\left[x_{s}+\frac{F_{s}\left(x_{s}\right)}{f_{s}\left(x_{s}\right)}\right]\right\} p\left(x_{b}, x_{s}\right) \\
& \quad f_{s} \mathrm{~d} x_{s} f_{b} \mathrm{~d} x_{b} \\
& \quad=\int_{c}^{d} \int_{c}^{x_{b} \wedge d}\left\{\left[x_{b}-\frac{1-F_{b}\left(x_{b}\right)}{f_{b}\left(x_{b}\right)}-x_{s}\right] f_{s}-F_{s}\left(x_{s}\right)\right\} f_{b} \mathrm{~d} x_{s} \mathrm{~d} x_{b}
\end{aligned}
$$

where the RHS uses the definition of $p\left(x_{b}, x_{s}\right)$.

$$
\begin{aligned}
= & \int_{c}^{d}\left[\left(x_{b}-\frac{1-F_{b}\left(x_{b}\right)}{f_{b}\left(x_{b}\right)}-x_{s}\right) F_{s}\left(x_{s}\right)\right]_{c}^{\left(x_{b} \wedge d\right)} f_{b} \mathrm{~d} x_{b} \\
= & \int_{c}^{d}\left[\left(x_{b}-\frac{1-F_{b}\left(x_{b}\right)}{f_{b}\left(x_{b}\right)}-\left(x_{b} \wedge d\right)\right) F_{s}\left(x_{b} \wedge d\right)\right] f_{b} \mathrm{~d} x_{b} \\
= & \int_{c}^{d}\left[-\left[1-F_{b}\left(x_{b}\right)\right] F_{s}\left(x_{b}\right)\right] \mathrm{d} x_{b} \\
& \quad+\int_{c}^{d}\left[\left(x_{b}-d\right) f_{b}\left(x_{b}\right)+\left(F_{b}\left(x_{b}\right)-1\right)\right] \mathrm{d} x_{b}
\end{aligned}
$$

QED

\section*{Alternative proof}
This proof is due to Williams (1999).

\begin{enumerate}
  \item Consider a DRM that satisfies BIC and ex post efficiency.
\end{enumerate}

\begin{itemize}
  \item If $x_{s}<x_{b}$ then
  \item Buyer pays $x_{s}$ to seller and gets object.
  \item Seller gives up object and receives $x_{b}$.
  \item If $x_{s}>x_{b}$ then
  \item there is no transaction and no transfers.
\end{itemize}

Note it is a DS to report valuation truthfully.\\
2. Corresponding expected transfers are

\begin{itemize}
  \item Buyer pays $T_{b}\left(x_{b}\right)=E_{x_{s}}\left[x_{s} \chi_{x_{s}<x_{b}}\left(x_{s}\right)\right]$.
  \item Seller receives $T_{s}\left(x_{s}\right)=E_{x_{b}}\left[x_{b} \chi_{x_{s}<x_{b}}\left(x_{b}\right)\right]$.
  \item Let $\left(p^{\prime}, t^{\prime}\right)$, be any mechanism that is ex post efficient. By the Expected Transfer Theorem, it must have the same expected transfers as the proposed mechanism up to a constant:
  \item Buyer pays $T_{b}^{\prime}\left(x_{b}\right)=E_{x_{s}}\left[x_{s} \chi_{x_{s}<x_{b}}\left(x_{s}\right)\right]+b$.
  \item Seller receives $T_{s}^{\prime}\left(x_{s}\right)=E_{x_{b}}\left[x_{b} \chi_{x_{s}<x_{b}}\left(x_{b}\right)\right]+s$.
\end{itemize}

\begin{enumerate}
  \setcounter{enumi}{2}
  \item Players expected payoffs in mechanism ( $p^{\prime}, t^{\prime}$ ).
\end{enumerate}

$$
\begin{aligned}
U_{b}^{\prime}\left(x_{b}\right) & =\overbrace{E_{x_{s}}\left[x_{b} \chi_{x_{s}<x_{b}}\right]}^{Q_{b}^{\prime}\left(x_{b}\right) \cdot x_{b}}-\overbrace{E_{x_{s}}\left[x_{s} \chi_{x_{s}<x_{b}}\left(x_{s}\right)\right]+b}^{T_{b}^{\prime}\left(x_{b}\right)} \\
& =E_{x_{s}}\left[\left(x_{b}-x_{s}\right) \chi_{x_{s}<x_{b}}\right]+b \\
U_{s}^{\prime}\left(x_{s}\right) & =T_{s}^{\prime}\left(x_{s}\right) \\
& =\overbrace{E_{x_{b}}\left[x_{b} \chi_{x_{s}<x_{b}}\right]+s}-\overbrace{E_{x_{b}}\left[x_{s} \chi_{x_{s}<x_{b}}\right]}^{Q_{s}^{\prime}\left(x_{s}\right) \cdot x_{s}} \\
& =E_{x_{b}}\left[\left(x_{b}-x_{s}\right) \chi_{x_{s}<x_{b}}\right]+s
\end{aligned}
$$

\begin{itemize}
  \item Using Individual Rationality we sign of $b$ and $s$ :
\end{itemize}

Recall that $x_{b}, x_{s} \in[c, d]$.\\
Then

\begin{itemize}
  \item $U_{b}^{\prime}(c) \geq 0 \Rightarrow b \geq 0$.
  \item $U_{s}^{\prime}(d) \geq 0 \Rightarrow s \geq 0$.
\end{itemize}

\begin{enumerate}
  \setcounter{enumi}{3}
  \item Total Ex Ante Payoff
\end{enumerate}

$$
\begin{aligned}
E\left[U_{b}^{\prime}\left(x_{b}\right)+U_{s}^{\prime}\left(x_{s}\right)\right] & =2 E\left[\left(x_{b}-x_{s}\right) \chi_{x_{s}<x_{b}}\right]+b+s \\
& >E\left[\left(x_{b}-x_{s}\right) \chi_{x_{s}<x_{b}}\right]
\end{aligned}
$$

Necessary payments to satisfy IR exceed the maximum total gains from trade (those prescribed by efficient outcome). QED

\section*{Remarks}
\begin{itemize}
  \item To induce seller and buyer to reveal their information, the mechanims must provide them certain payoffs.
  \item To satisfy IR constraints, those payoffs are high because IR requires a minimum payoff for certain types.
  \item An expost efficient mechanism cannot get enough expected revenue to satisfy the IR constraints.
  \item Balance budget is not possible.
\end{itemize}

\section*{References}
Arrow, Kenneth J. 1979. "The Property Rights Doctrine and Demand Revelation Under Incomplete Information." In Economics and Human Welfare, by M. Boskin. New York: Academic Press. d'Aspremont, Claude, and Louis-André Gérard-Varet. 1979. "Incentives and Incomplete Information." Journal of Public Economics 11: 25-45.

Myerson, Roger. 1981. "Optimal Auction Design." Mathematics of Operations Research 6: 58-73.\\
Myerson, Roger, and Mark Satterthwaite. 1983. "Efficient Mechanisms for Bilateral Trading." Journal of Economic Theory 29: 265-81.\\
Vickrey, W. 1961. "Counterspeculation, Auctions and Competitive Sealed Tenders." Journal of Finance 16: 8-37.

Williams, Steven. 1999. "A Characterization of Efficient, Bayesian Incentive Compatible Mechanisms." Economic Theory 14: 155-80.


\end{document}